% ==========================================================================
%  Chapter 77 -- Multiscale Architecture Refactoring: Implementation
% ==========================================================================

\chapter{Multiscale Architecture Refactoring: Implementation}
\label{ch:multiscale-refactoring}


% =========================================================================
\section{Introduction}
\label{sec:refactor-intro}

Chapter~\ref{ch:multiscale-architecture} identified six rigidities in
the FE\textsuperscript{2} coupling implementation and proposed a
nine-phase refactoring plan to achieve a SOLID-conformant,
extensible, and HPC-ready architecture.  Chapter~\ref{ch:build-system-optimisation}
diagnosed the root causes of the excessive compilation time and
proposed a mitigation plan.  This chapter documents the
implementation of the first seven phases of that plan:

\begin{enumerate}[label=\textbf{Phase \arabic*}:]
  \setcounter{enumi}{-1}
  \item Build system cleanup: consolidation of 20 duplicate seismic
    targets.
  \item \texttt{LocalModelAdapter} concept and type-erased handle.
  \item \texttt{MaterialFactory} injection into the sub-model solver.
  \item \texttt{MultiscaleModel} aggregation class.
  \item \texttt{CouplingStrategy} interfaces and \texttt{MultiscaleAnalysis}
    orchestrator.
  \item \texttt{HomogenisationStrategy} for switchable upscaling.
\end{enumerate}

All code changes presented here compile under GCC~15.2 with
\texttt{-Werror} enabled.  The full test suite (63~CTest cases)
continues to pass after each phase.

Section~\ref{sec:refactor-build-cleanup} covers the build system
cleanup.  Sections~\ref{sec:refactor-local-model}
through~\ref{sec:refactor-homogenisation} present each architectural
phase with design rationale, code listings, and SOLID analysis.
Section~\ref{sec:refactor-orchestrator} presents the
\texttt{MultiscaleAnalysis} orchestrator that absorbs the staggered
coupling loop.  Section~\ref{sec:refactor-discussion} provides a
comprehensive discussion of design decisions and trade-offs.


% =========================================================================
\section{Phase 0.1: Build System Cleanup}
\label{sec:refactor-build-cleanup}

% -------------------------------------------------------------------------
\subsection{Problem: 20 Duplicate Seismic Targets}

As diagnosed in Chapter~\ref{ch:build-system-optimisation}, the
CMake build system contained 20 near-identical executable targets
(\texttt{fall\_n\_seismic\_v2} through \texttt{fall\_n\_seismic\_v21}),
each compiling the same source file
(\texttt{main\_multiscale\_seismic.cpp}) with a different preprocessor
definition.  Each target incurred the full parse cost of
\texttt{header\_files.hh} ($\sim$100 project headers), resulting in
${\sim}35\%$ of the total build time spent on redundant compilation.

\subsection{Solution: Target Consolidation}

All 20 duplicate target blocks were removed from
\texttt{CMakeLists.txt}.  The single pre-existing target
\texttt{fall\_n\_seismic} is retained; variant selection is
performed at runtime via command-line arguments.

\paragraph{Impact.}
\begin{itemize}
  \item Ninja targets: $115 \to 88$ ($-23\%$).
  \item Build time: estimated ${\sim}35\%$ reduction for full rebuilds.
  \item No test regressions: 63/63 CTest cases pass.
\end{itemize}


% =========================================================================
\section{Phase 1: \texttt{LocalModelAdapter} Concept}
\label{sec:refactor-local-model}

% -------------------------------------------------------------------------
\subsection{Design Rationale}
\label{sec:refactor-local-model-rationale}

The \texttt{NonlinearSubModelEvolver} class
(Chapter~\ref{ch:nl-sub-model-evolver}) is the sole concrete
sub-model solver.  However, the multiscale analysis framework should
be open to alternative sub-model types (e.g.\ a future
\texttt{LinearSubModel} for elastic columns, or a
\texttt{ReducedOrderModel} for computational efficiency).

The design choice between a C++20 concept and a classical abstract
base class (ABC) was guided by consistency with the existing
codebase:

\begin{itemize}
  \item \texttt{FiniteElement}
    (\texttt{src/elements/FiniteElementConcept.hh}) uses a concept.
  \item \texttt{SteppableSolver}
    (\texttt{src/analysis/SteppableSolverConcept.hh}) uses a concept.
  \item \texttt{IncrementalControlPolicy}
    (\texttt{src/analysis/IncrementalControl.hh}) uses a concept.
  \item \texttt{DamageCriterion}
    (\texttt{src/analysis/DamageCriterion.hh}) uses an ABC.
\end{itemize}

Concepts are preferred for interfaces where compile-time dispatch
suffices.  Since the inner staggered loop processes sub-models of
homogeneous type (all are \texttt{NonlinearSubModelEvolver}), a
concept avoids the overhead of virtual dispatch.  For the
heterogeneous case (mixing different solver types in one container),
a type-erased handle is provided.

% -------------------------------------------------------------------------
\subsection{The \texttt{LocalModelAdapter} Concept}
\label{sec:refactor-concept-def}

The concept is defined in
\texttt{src/reconstruction/LocalModelAdapter.hh}:

\begin{lstlisting}[caption={C++20 concept for local sub-model solvers.},
                   label={lst:local-model-concept}]
template <typename T>
concept LocalModelAdapter = requires(
    T& t,
    const SectionKinematics& kin,
    double width, double height, double h_pert, double time)
{
    { t.update_kinematics(kin, kin) };
    { t.solve_step(time) };
    { t.section_tangent(width, height, h_pert) }
        -> std::convertible_to<Eigen::Matrix<double, 6, 6>>;
    { t.section_forces(width, height) }
        -> std::convertible_to<Eigen::Vector<double, 6>>;
    { t.commit_state() };
    { t.revert_state() };
    { t.parent_element_id() }
        -> std::convertible_to<std::size_t>;
};
\end{lstlisting}

Each required expression maps to a well-defined responsibility:
\begin{description}
  \item[\texttt{update\_kinematics}] Localisation: imposes macro-scale
    beam kinematics as Dirichlet BCs on the sub-model.
  \item[\texttt{solve\_step}] Solves the nonlinear boundary-value
    problem from the current converged state.
  \item[\texttt{section\_tangent}] Computational homogenisation:
    returns the $6\times6$ tangent $\mathbf{D}_{\mathrm{hom}} =
    \partial\mathbf{s}/\partial\mathbf{e}$ via perturbation.
  \item[\texttt{section\_forces}] Returns the upscaled section forces
    $\mathbf{s} = [N, M_y, M_z, V_y, V_z, M_t]^\top$.
  \item[\texttt{commit\_state} / \texttt{revert\_state}]
    Transaction semantics for irreversible constitutive state
    (damage, plasticity).
  \item[\texttt{parent\_element\_id}] Identifies the owning beam
    element in the global model.
\end{description}


% -------------------------------------------------------------------------
\subsection{Type-Erased Handle: \texttt{LocalModelHandle}}
\label{sec:refactor-type-erasure}

For heterogeneous containers (\texttt{std::vector<LocalModelHandle>}),
a type-erased wrapper follows the Concept/Model/Handle pattern
established by \texttt{FEM\_Element}
(\texttt{src/elements/FEM\_Element.hh}):

\begin{lstlisting}[caption={Type-erased handle (abridged).},
                   label={lst:local-model-handle}]
class LocalModelHandle {
    struct Concept {  // pure virtual interface
        virtual ~Concept() = default;
        virtual void update_kinematics(...) = 0;
        virtual void solve_step(double) = 0;
        virtual Eigen::Matrix<double,6,6>
            section_tangent(double, double, double) = 0;
        virtual Eigen::Vector<double,6>
            section_forces(double, double) = 0;
        virtual void commit_state() = 0;
        virtual void revert_state() = 0;
        virtual std::size_t parent_element_id() const = 0;
    };

    template <LocalModelAdapter T>
    struct Model final : Concept {
        T adapter_;
        // ... delegates all calls to adapter_
    };

    std::unique_ptr<Concept> pimpl_;
public:
    template <LocalModelAdapter T>
    LocalModelHandle(T adapter);  // implicit conversion

    // Move-only (sub-models own PETSc resources)
    LocalModelHandle(LocalModelHandle&&) noexcept = default;
    LocalModelHandle(const LocalModelHandle&) = delete;

    // Forwarding interface mirrors the concept
    void update_kinematics(const SectionKinematics&,
                           const SectionKinematics&);
    void solve_step(double);
    auto section_tangent(double, double, double = 1e-6)
        -> Eigen::Matrix<double,6,6>;
    // ...
};
\end{lstlisting}

The handle is \emph{move-only} because sub-model solvers own PETSc
objects (DM, SNES, Vec) that cannot be trivially copied.  The
implicit converting constructor allows any
\texttt{LocalModelAdapter}-satisfying type to be stored in a
\texttt{std::vector<LocalModelHandle>} without explicit wrapping.

% -------------------------------------------------------------------------
\subsection{Adapting \texttt{NonlinearSubModelEvolver}}
\label{sec:refactor-evolver-adaptation}

The evolver was modified to satisfy the concept:

\begin{enumerate}
  \item \textbf{Public aliases:} New methods \texttt{section\_tangent()}
    and \texttt{section\_forces()} delegate to the existing private
    methods \texttt{compute\_homogenized\_tangent()} and
    \texttt{compute\_homogenized\_forces()}, preserving backward
    compatibility:
\begin{lstlisting}[caption={Public aliases on the evolver.},
                   label={lst:evolver-aliases}]
// LocalModelAdapter concept conformance
[[nodiscard]] Eigen::Matrix<double, 6, 6>
section_tangent(double width, double height,
                double h_pert = 1.0e-6) {
    return compute_homogenized_tangent(width, height, h_pert);
}

[[nodiscard]] Eigen::Vector<double, 6>
section_forces(double width, double height) {
    return compute_homogenized_forces(width, height);
}
\end{lstlisting}

  \item \textbf{Public state management:} Methods
    \texttt{commit\_state()} and \texttt{revert\_state()} were moved
    from the private section to public.  \texttt{commit\_state()}
    traverses the DM mesh and calls \texttt{commit\_material\_state()}
    on each element.  \texttt{revert\_state()} calls
    \texttt{revert\_material\_state()} to discard the trial state
    after a perturbation solve.

  \item \textbf{Static assertion:} A compile-time check at the end
    of the header ensures concept satisfaction is maintained as the
    class evolves:
\begin{lstlisting}[caption={Compile-time concept check.}]
static_assert(
    LocalModelAdapter<NonlinearSubModelEvolver>,
    "NonlinearSubModelEvolver must satisfy "
    "the LocalModelAdapter concept");
\end{lstlisting}
\end{enumerate}


% =========================================================================
\section{Phase 2: \texttt{MaterialFactory} Injection}
\label{sec:refactor-material-factory}

% -------------------------------------------------------------------------
\subsection{Problem: Hardcoded Constitutive Models}

The \texttt{first\_solve()} method of
\texttt{NonlinearSubModelEvolver} directly instantiated
\texttt{KoBatheConcrete3D} and \texttt{MenegottoPintoSteel} material
models:

\begin{lstlisting}[caption={Hardcoded material construction (before).}]
// In first_solve():
InelasticMaterial<KoBatheConcrete3D> mat_inst{fc_};
Material<ThreeDimensionalMaterial> mat{mat_inst, InelasticUpdate{}};
// ...
MenegottoPintoSteel steel_impl{200000.0, 420.0, 0.01};
InelasticMaterial<MenegottoPintoSteel> steel_inst{std::move(steel_impl)};
Material<UniaxialMaterial> steel_mat{std::move(steel_inst), InelasticUpdate{}};
\end{lstlisting}

This violated the Open--Closed Principle (OCP): using a different
concrete model (e.g.\ \texttt{Mazars3D}) or rebar steel (e.g.\
\texttt{BilinearSteel}) required modifying the evolver source.

% -------------------------------------------------------------------------
\subsection{Solution: Abstract Factory Pattern}
\label{sec:refactor-factory-design}

Two abstract factory interfaces are defined in
\texttt{src/reconstruction/MaterialFactory.hh}, following the
\texttt{DamageCriterion} ABC pattern from
\texttt{src/analysis/DamageCriterion.hh}:

\begin{lstlisting}[caption={Material factory ABCs.},
                   label={lst:material-factories}]
struct ConcreteMaterialFactory {
    virtual ~ConcreteMaterialFactory() = default;
    virtual Material<ThreeDimensionalMaterial> create() const = 0;
    virtual std::unique_ptr<ConcreteMaterialFactory> clone() const = 0;
};

struct RebarMaterialFactory {
    virtual ~RebarMaterialFactory() = default;
    virtual Material<UniaxialMaterial> create() const = 0;
    virtual std::unique_ptr<RebarMaterialFactory> clone() const = 0;
};
\end{lstlisting}

Default implementations encapsulate the previously hardcoded
construction:

\begin{lstlisting}[caption={Default factory implementations.},
                   label={lst:default-factories}]
class KoBatheConcreteMaterialFactory final
    : public ConcreteMaterialFactory
{
    double fc_;
public:
    explicit KoBatheConcreteMaterialFactory(double fc_MPa)
        : fc_{fc_MPa} {}

    Material<ThreeDimensionalMaterial> create() const override {
        InelasticMaterial<KoBatheConcrete3D> mat_inst{fc_};
        return Material<ThreeDimensionalMaterial>{
            mat_inst, InelasticUpdate{}};
    }

    std::unique_ptr<ConcreteMaterialFactory>
    clone() const override {
        return std::make_unique<
            KoBatheConcreteMaterialFactory>(*this);
    }
};
\end{lstlisting}

% -------------------------------------------------------------------------
\subsection{Injection into the Evolver}
\label{sec:refactor-factory-injection}

The evolver now holds two factory members:

\begin{lstlisting}[caption={Factory members in the evolver.}]
std::unique_ptr<ConcreteMaterialFactory> concrete_factory_;
std::unique_ptr<RebarMaterialFactory>    rebar_factory_;
\end{lstlisting}

Two constructors are provided:

\begin{enumerate}
  \item \textbf{Original constructor} (backward-compatible): initialises
    default factories using \texttt{KoBatheConcreteMaterialFactory}
    and \texttt{MenegottoPintoRebarFactory} with the same parameters
    as the previously hardcoded construction.
  \item \textbf{Factory-injecting constructor}: accepts custom
    \texttt{std::unique\_ptr<ConcreteMaterialFactory>} and
    \texttt{std::unique\_ptr<RebarMaterialFactory>}, enabling
    user-defined materials without modifying the evolver.
\end{enumerate}

In \texttt{first\_solve()}, the hardcoded material lines are replaced
by factory calls:

\begin{lstlisting}[caption={Factory-delegated material creation.}]
// Before: InelasticMaterial<KoBatheConcrete3D> mat_inst{fc_};
Material<Policy> mat = concrete_factory_->create();

// Before: MenegottoPintoSteel steel_impl{...};
Material<UniaxialMaterial> steel_mat = rebar_factory_->create();
\end{lstlisting}

\paragraph{OCP compliance.}
Adding a new constitutive model (e.g.\ \texttt{Mazars3D}) now
requires only writing a new factory subclass; the evolver source is
\emph{closed for modification}.


% =========================================================================
\section{Phase 3: \texttt{MultiscaleModel} Aggregation}
\label{sec:refactor-multiscale-model}

% -------------------------------------------------------------------------
\subsection{Problem: Scattered Model--Sub-Model Coupling}

In the original \texttt{main\_lshaped\_multiscale.cpp}, the
relationship between the global beam model and the local sub-model
solvers was managed through three separate variables:

\begin{lstlisting}[caption={Scattered state in original main.}]
std::vector<NonlinearSubModelEvolver> nl_evolvers;  // sub-models
std::vector<Eigen::Matrix<double,6,6>> D_prev;      // tangent history
// kinematics extraction: a lambda capturing model and material params
auto extract_beam_kinematics = [&](std::size_t e_idx) { ... };
\end{lstlisting}

The kinematics extraction lambda captured the global model, element
type, and material constants (\texttt{EC\_COL}, \texttt{GC\_COL},
\texttt{NU\_RC}) by reference.  The injection of homogenised
tangent and forces was performed inline via
\texttt{beam->set\_homogenized\_tangent(D\_hom)}.  This scattered
organisation made it impossible to reuse the coupling logic in
another context (different building, different element type).

% -------------------------------------------------------------------------
\subsection{Design: Template on Local Model Type}
\label{sec:refactor-model-design}

The \texttt{MultiscaleModel} class aggregates the global--local
relationship into a single reusable object, defined in
\texttt{src/analysis/MultiscaleModel.hh}:

\begin{lstlisting}[caption={\texttt{MultiscaleModel} class template (abridged).},
                   label={lst:multiscale-model}]
template <LocalModelAdapter LocalModel>
class MultiscaleModel {

    std::vector<LocalModel> local_models_;
    std::unordered_map<std::size_t, std::size_t> transition_map_;

    std::function<ElementKinematics(std::size_t)>
        kinematics_extractor_;

    using InjectorFn = std::function<void(
        std::size_t,
        const Eigen::Matrix<double,6,6>&,
        const Eigen::Vector<double,6>&)>;
    InjectorFn response_injector_;

public:
    void register_local_model(std::size_t parent_elem_id,
                              LocalModel model);

    LocalModel& local_model_for(std::size_t elem_id);

    ElementKinematics extract_kinematics(std::size_t elem_id) const;

    void inject_response(std::size_t elem_id,
                         const Eigen::Matrix<double,6,6>& D_hom,
                         const Eigen::Vector<double,6>& f_hom);

    std::size_t num_local_models() const;
    std::vector<std::size_t> parent_element_ids() const;
};
\end{lstlisting}

Key design decisions:

\begin{enumerate}
  \item \textbf{Template on local model type.}  The class is
    parametrised on a \texttt{LocalModelAdapter}-constrained type
    rather than using \texttt{LocalModelHandle} directly.  This
    avoids virtual dispatch overhead in the inner staggered loop,
    while the concept constraint ensures type safety.  If
    heterogeneous sub-models are needed, instantiate with
    \texttt{LocalModelHandle}.

  \item \textbf{Callback-based scale bridging.}  The kinematics
    extraction and response injection functions are stored as
    \texttt{std::function} callbacks, set via
    \texttt{set\_kinematics\_extractor()} and
    \texttt{set\_response\_injector()}.  This decouples the model
    aggregation from the concrete element type
    (\texttt{TimoshenkoBeam3D}, \texttt{EulerBernoulliBeam}, etc.)
    and the specific model template parameters.

  \item \textbf{Transition map.}  An \texttt{unordered\_map} maps
    global beam element IDs to sub-model indices, enabling $O(1)$
    lookup.  This replaces the implicit index-based correspondence
    in the original code.
\end{enumerate}


% =========================================================================
\section{Phase 4: Coupling Strategy Interfaces}
\label{sec:refactor-coupling-strategy}

% -------------------------------------------------------------------------
\subsection{Problem: Hardcoded Coupling Parameters}

The staggered coupling loop in
\texttt{main\_lshaped\_multiscale.cpp} used four hardcoded
compile-time constants:

\begin{lstlisting}[caption={Hardcoded coupling parameters.}]
static constexpr int    MAX_STAGGERED_ITER  = 4;
static constexpr double STAGGERED_TOL       = 0.05;
static constexpr double STAGGERED_RELAX     = 0.7;
static constexpr int    COUPLING_START_STEP = 10;
\end{lstlisting}

The convergence criterion (relative Frobenius norm), relaxation
policy (constant $\omega$), and coupling mode (two-way staggered)
were all embedded in the loop body.  Switching to one-way coupling
or Aitken adaptive relaxation required editing the main driver.

% -------------------------------------------------------------------------
\subsection{Strategy Interfaces}
\label{sec:refactor-strategy-interfaces}

Three abstract base classes are defined in
\texttt{src/analysis/CouplingStrategy.hh}, implementing the
\emph{Strategy pattern} \cite{GammaDesignPatterns}:

\paragraph{ScaleBridgePolicy.}
Controls whether the coupling is one-way (sub-models are
\emph{passive observers}) or two-way (sub-models feed back
stiffness and forces to the global model):

\begin{lstlisting}[caption={Scale bridge policy ABC and concrete implementations.},
                   label={lst:scale-bridge}]
struct ScaleBridgePolicy {
    virtual ~ScaleBridgePolicy() = default;
    virtual bool requires_feedback() const = 0;
    virtual int  max_staggered_iterations() const = 0;
};

class OneWayCoupling final : public ScaleBridgePolicy {
public:
    bool requires_feedback() const override { return false; }
    int  max_staggered_iterations() const override { return 1; }
};

class TwoWayStaggered final : public ScaleBridgePolicy {
    int max_iter_;
public:
    explicit TwoWayStaggered(int max_iter = 4)
        : max_iter_{max_iter} {}
    bool requires_feedback() const override { return true; }
    int  max_staggered_iterations() const override
        { return max_iter_; }
};
\end{lstlisting}

\paragraph{CouplingConvergence.}
Defines the convergence criterion for the staggered iteration.
The current implementation uses the relative Frobenius norm:

\begin{lstlisting}[caption={Frobenius convergence criterion.},
                   label={lst:frobenius-conv}]
struct CouplingConvergence {
    virtual ~CouplingConvergence() = default;
    virtual bool converged(
        std::span<const Eigen::Matrix<double,6,6>> D_prev,
        std::span<const Eigen::Matrix<double,6,6>> D_curr) const = 0;
};

class FrobeniusConvergence final : public CouplingConvergence {
    double tol_;
public:
    explicit FrobeniusConvergence(double tol = 0.05)
        : tol_{tol} {}
    bool converged(
        std::span<const Eigen::Matrix<double,6,6>> D_prev,
        std::span<const Eigen::Matrix<double,6,6>> D_curr
    ) const override {
        for (std::size_t i = 0; i < D_prev.size(); ++i) {
            double d_norm = D_curr[i].norm();
            double delta  = (D_curr[i] - D_prev[i]).norm();
            if (d_norm > 1e-14 && delta / d_norm > tol_)
                return false;
        }
        return true;
    }
};
\end{lstlisting}

The convergence criterion checks whether the relative change in \emph{every}
sub-model's tangent matrix is below the tolerance~$\tau$:
\begin{equation}
  \frac{\bigl\lVert \mathbf{D}_{\mathrm{hom},i}^{(k)}
        - \mathbf{D}_{\mathrm{hom},i}^{(k-1)} \bigr\rVert_F}
       {\bigl\lVert \mathbf{D}_{\mathrm{hom},i}^{(k)} \bigr\rVert_F}
  < \tau, \qquad \forall i \in \{1,\ldots,N_s\}
  \label{eq:frobenius-conv}
\end{equation}
where $N_s$ is the number of sub-models and $k$ is the staggered
iteration counter.

\paragraph{RelaxationPolicy.}
Controls how the updated tangent is blended with the previous
iteration's value:

\begin{lstlisting}[caption={Relaxation policy implementations.},
                   label={lst:relaxation}]
struct RelaxationPolicy {
    virtual ~RelaxationPolicy() = default;
    virtual void relax(
        Eigen::Matrix<double,6,6>& D_new,
        const Eigen::Matrix<double,6,6>& D_prev,
        int iteration) = 0;
};

class ConstantRelaxation final : public RelaxationPolicy {
    double omega_;
public:
    explicit ConstantRelaxation(double omega = 0.7)
        : omega_{omega} {}
    void relax(
        Eigen::Matrix<double,6,6>& D_new,
        const Eigen::Matrix<double,6,6>& D_prev,
        [[maybe_unused]] int iteration) override
    {
        D_new = omega_ * D_new + (1.0 - omega_) * D_prev;
    }
};
\end{lstlisting}

The constant relaxation applies:
\begin{equation}
  \mathbf{D}_{\mathrm{hom}}^{(k)}
    \leftarrow
    \omega \, \mathbf{D}_{\mathrm{hom}}^{(k)}
    + (1-\omega) \, \mathbf{D}_{\mathrm{hom}}^{(k-1)}
  \label{eq:relaxation}
\end{equation}
with $\omega = 0.7$ as default.  A \texttt{NoRelaxation} policy is
also provided for cases where relaxation is unnecessary or
destabilising.


% =========================================================================
\section{Phase 4.3: \texttt{MultiscaleAnalysis} Orchestrator}
\label{sec:refactor-orchestrator}

% -------------------------------------------------------------------------
\subsection{Design}

The \texttt{MultiscaleAnalysis} class, defined in
\texttt{src/analysis/MultiscaleAnalysis.hh}, absorbs the
staggered coupling loop into a reusable, strategy-controlled
template:

\begin{lstlisting}[caption={\texttt{MultiscaleAnalysis} class (abridged).},
                   label={lst:multiscale-analysis}]
template <LocalModelAdapter LocalModel>
class MultiscaleAnalysis {
    MultiscaleModel<LocalModel>             model_;
    std::unique_ptr<ScaleBridgePolicy>      bridge_;
    std::unique_ptr<CouplingConvergence>    convergence_;
    std::unique_ptr<RelaxationPolicy>       relaxation_;

    int    coupling_start_step_{0};
    double section_width_{0.30};
    double section_height_{0.30};
    double tangent_perturbation_{1.0e-6};

    std::vector<Eigen::Matrix<double,6,6>> D_prev_;

public:
    MultiscaleAnalysis(
        MultiscaleModel<LocalModel> model,
        std::unique_ptr<ScaleBridgePolicy> bridge,
        std::unique_ptr<CouplingConvergence> convergence,
        std::unique_ptr<RelaxationPolicy> relaxation);

    void set_coupling_start_step(int step);
    void set_section_dimensions(double w, double h);

    // One global time step + staggered loop
    bool step(double time, int global_step);
};
\end{lstlisting}

% -------------------------------------------------------------------------
\subsection{The \texttt{step()} Method}
\label{sec:refactor-step-method}

The \texttt{step()} method implements the full staggered coupling
algorithm for one global time step.  The pseudocode is:

\begin{enumerate}
  \item Determine whether coupling is active:
    \texttt{bridge\_->requires\_feedback()} \textbf{and}
    \texttt{global\_step $\geq$ coupling\_start\_step\_}.
  \item For each staggered iteration $k = 0, 1, \ldots$:
    \begin{enumerate}[label=(\alph*)]
      \item \textbf{Downscale.}  For each local model:
        extract kinematics from the global model via
        \texttt{model\_.extract\_kinematics()}, update the sub-model
        BCs, and solve the nonlinear BVP.
      \item \textbf{Upscale.}  For each local model:
        compute the homogenised tangent
        $\mathbf{D}_{\mathrm{hom}}^{(k)}$ via perturbation,
        apply relaxation $(\ref{eq:relaxation})$, and check
        convergence $(\ref{eq:frobenius-conv})$.
      \item \textbf{Inject.}  Write
        $\mathbf{D}_{\mathrm{hom}}^{(k)}$ and
        $\mathbf{f}_{\mathrm{hom}}^{(k)}$ into the global beam
        elements via
        \texttt{model\_.inject\_response()}.
      \item If $k > 0$ and converged, \textbf{break}.
    \end{enumerate}
  \item Commit the converged constitutive state of all sub-models.
\end{enumerate}

\paragraph{Intended usage in the main driver:}
\begin{lstlisting}[caption={Proposed refactored main driver.},
                   label={lst:proposed-main}]
auto bridge = std::make_unique<TwoWayStaggered>(4);
auto conv   = std::make_unique<FrobeniusConvergence>(0.05);
auto relax  = std::make_unique<ConstantRelaxation>(0.7);

MultiscaleAnalysis<NonlinearSubModelEvolver> analysis(
    std::move(ms_model),
    std::move(bridge), std::move(conv), std::move(relax));
analysis.set_coupling_start_step(10);
analysis.set_section_dimensions(COL_B, COL_H);

for (int step = 0; step < n_steps; ++step)
    analysis.step(time[step], step);
\end{lstlisting}

This replaces the ${\sim}80$ lines of staggered loop logic in the
original \texttt{main\_lshaped\_multiscale.cpp} with a single
\texttt{analysis.step()} call, while retaining full configurability
via the injected strategy objects.


% =========================================================================
\section{Phase 5: \texttt{HomogenisationStrategy}}
\label{sec:refactor-homogenisation}

% -------------------------------------------------------------------------
\subsection{Problem: Two Upscaling Methods}

The sub-model solver implements two distinct methods for extracting
section forces from the 3-D displacement field:

\begin{enumerate}
  \item \textbf{Boundary reaction} (used by
    \texttt{section\_forces\_from\_reactions()}):
    assembles the internal force vector $\mathbf{f}_{\mathrm{int}}$,
    reads the nodal forces at Face~B nodes, and computes the section
    resultants $\mathbf{s} = [N, M_y, M_z, V_y, V_z, M_t]^\top$
    from the reaction force and moment sums.  This method is
    consistent with the perturbation tangent computation and gives
    equilibrium-consistent forces.

  \item \textbf{Volume averaging} (used by
    \texttt{compute\_homogenized\_forces()} via the
    \texttt{homogenize()} free function in
    \texttt{HomogenizedSection.hh}):
    integrates the Cauchy stress field over the cross-section using
    the uniform-stress approximation.  This is the classical
    Hill--Mandel approach.
\end{enumerate}

The choice between these methods was previously hardcoded: the
tangent used boundary reactions while the forces used volume
averaging.

% -------------------------------------------------------------------------
\subsection{Strategy Interface}

The \texttt{HomogenisationStrategy} ABC in
\texttt{src/reconstruction/HomogenisationStrategy.hh} provides a
clean selection mechanism:

\begin{lstlisting}[caption={Homogenisation strategy interface.},
                   label={lst:homogenisation-strategy}]
struct HomogenisationStrategy {
    virtual ~HomogenisationStrategy() = default;

    virtual Eigen::Vector<double, 6> select_forces(
        const Eigen::Vector<double, 6>& boundary_reaction,
        const Eigen::Vector<double, 6>& volume_average
    ) const = 0;

    virtual std::string name() const = 0;
};
\end{lstlisting}

The interface is \emph{deliberately independent} of the solver class
to avoid circular header dependencies.  The sub-model solver
computes both force vectors and passes them to the strategy for
selection.

\paragraph{Concrete implementations:}

\begin{lstlisting}[caption={Concrete homogenisation strategies.},
                   label={lst:homogenisation-impls}]
struct BoundaryReactionHomogenisation final
    : HomogenisationStrategy
{
    Eigen::Vector<double, 6> select_forces(
        const Eigen::Vector<double, 6>& boundary_reaction,
        const Eigen::Vector<double, 6>& /*volume_average*/
    ) const override { return boundary_reaction; }

    std::string name() const override
        { return "boundary-reaction"; }
};

struct VolumeAverageHomogenisation final
    : HomogenisationStrategy
{
    Eigen::Vector<double, 6> select_forces(
        const Eigen::Vector<double, 6>& /*boundary_reaction*/,
        const Eigen::Vector<double, 6>& volume_average
    ) const override { return volume_average; }

    std::string name() const override
        { return "volume-average"; }
};
\end{lstlisting}

The \texttt{BoundaryReactionHomogenisation} strategy is recommended
as the default because it produces forces consistent with the
perturbation-based tangent computation (both use assembled
$\mathbf{f}_{\mathrm{int}}$), whereas the volume-average approach
may give stale results when the stored material state is behind the
displacement field.


% =========================================================================
\section{Discussion}
\label{sec:refactor-discussion}

% -------------------------------------------------------------------------
\subsection{SOLID Compliance}
\label{sec:refactor-solid}

Table~\ref{tab:solid-compliance} summarises how each refactoring
phase addresses specific SOLID principles.

\begin{table}[htbp]
  \centering
  \caption{SOLID compliance of each refactoring phase.}
  \label{tab:solid-compliance}
  \small
  \begin{tabular}{@{}llp{7cm}@{}}
    \toprule
    \textbf{Phase} & \textbf{Principle} & \textbf{Mechanism} \\
    \midrule
    1 -- LocalModelAdapter &
      OCP, DIP &
      New sub-model types can be added without modifying the
      orchestrator; high-level modules depend on the concept
      abstraction. \\[4pt]
    2 -- MaterialFactory &
      OCP, SRP &
      New materials via factory subclasses; material creation
      responsibility extracted from the solver. \\[4pt]
    3 -- MultiscaleModel &
      SRP &
      Model--sub-model relationship encapsulated; main driver cleaned. \\[4pt]
    4 -- CouplingStrategy &
      OCP, ISP, SRP &
      Convergence, relaxation, and coupling mode are independently
      switchable; each interface has a single, narrow contract. \\[4pt]
    5 -- HomogenisationStrategy &
      OCP &
      Upscaling method switchable without modifying the solver. \\
    \bottomrule
  \end{tabular}
\end{table}

% -------------------------------------------------------------------------
\subsection{Concept vs.\ ABC Trade-Offs}
\label{sec:refactor-concept-vs-abc}

The choice between C++20 concepts and ABCs was made on a
case-by-case basis:

\begin{itemize}
  \item \textbf{Concept} (\texttt{LocalModelAdapter}): The local
    model type is known at compile time for a given analysis.  A
    concept avoids virtual dispatch in the inner loop (perturbation
    tangent: 6 SNES solves per sub-model per staggered iteration)
    and enables the compiler to inline critical paths.  The
    type-erased handle \texttt{LocalModelHandle} is available as an
    opt-in for the heterogeneous case.

  \item \textbf{ABC} (\texttt{ConcreteMaterialFactory},
    \texttt{ScaleBridgePolicy}, \texttt{CouplingConvergence},
    \texttt{RelaxationPolicy}, \texttt{HomogenisationStrategy}):
    These are called once per staggered iteration (not in a tight
    loop), so the virtual dispatch cost is negligible.  ABCs provide
    runtime polymorphism, enabling strategy selection at
    configuration time rather than compilation time.
\end{itemize}

% -------------------------------------------------------------------------
\subsection{Circular Dependency Avoidance}
\label{sec:refactor-circular-deps}

The \texttt{HomogenisationStrategy} interface was redesigned to
avoid a circular dependency between
\texttt{HomogenisationStrategy.hh} and
\texttt{NonlinearSubModelEvolver.hh}.  The original design took a
\texttt{NonlinearSubModelEvolver\&} parameter.  The final design
takes pre-computed force vectors (boundary reaction and volume
average), inverting the dependency: the solver computes both and
passes them to the strategy.  This also makes the strategies
testable in isolation.

% -------------------------------------------------------------------------
\subsection{Type-Erasure Pattern Reuse}
\label{sec:refactor-type-erasure-pattern}

The \texttt{LocalModelHandle} follows the same Concept/Model/Handle
triple used by \texttt{FEM\_Element}
(\texttt{src/elements/FEM\_Element.hh}).  This pattern is now a
well-established idiom in the codebase:

\begin{center}
\begin{tabular}{lll}
  \toprule
  \textbf{Handle} & \textbf{Concept} & \textbf{Purpose} \\
  \midrule
  \texttt{FEM\_Element} & \texttt{FiniteElement} & Heterogeneous element containers \\
  \texttt{LocalModelHandle} & \texttt{LocalModelAdapter} & Heterogeneous sub-model containers \\
  \bottomrule
\end{tabular}
\end{center}

Both are move-only (owning PETSc/mesh resources) and use
\texttt{std::unique\_ptr<Concept>} for type erasure.

% -------------------------------------------------------------------------
\subsection{Build Verification}
\label{sec:refactor-build-verification}

After all phases:

\begin{itemize}
  \item \textbf{Compilation:} 88 Ninja targets (down from 115),
    all compile under GCC~15.2 with \texttt{-Werror -std=c++23}.
  \item \textbf{Tests:} 63/63 CTest cases pass (51~unit, 6~integration,
    6~slow).
  \item \textbf{New headers:} 6 files added, listed in
    Table~\ref{tab:new-files}.
  \item \textbf{Modified headers:} 3 files modified, listed in
    Table~\ref{tab:modified-files}.
\end{itemize}

\begin{table}[htbp]
  \centering
  \caption{New files introduced by the refactoring.}
  \label{tab:new-files}
  \small
  \begin{tabular}{@{}ll@{}}
    \toprule
    \textbf{File} & \textbf{Purpose} \\
    \midrule
    \texttt{src/reconstruction/LocalModelAdapter.hh} &
      Concept + type-erased handle \\
    \texttt{src/reconstruction/MaterialFactory.hh} &
      Material factory ABCs + defaults \\
    \texttt{src/reconstruction/HomogenisationStrategy.hh} &
      Upscaling strategy ABC + implementations \\
    \texttt{src/analysis/CouplingStrategy.hh} &
      Coupling strategy ABCs + implementations \\
    \texttt{src/analysis/MultiscaleModel.hh} &
      Model aggregation template \\
    \texttt{src/analysis/MultiscaleAnalysis.hh} &
      Staggered coupling orchestrator \\
    \bottomrule
  \end{tabular}
\end{table}

\begin{table}[htbp]
  \centering
  \caption{Modified files.}
  \label{tab:modified-files}
  \small
  \begin{tabular}{@{}lp{8cm}@{}}
    \toprule
    \textbf{File} & \textbf{Changes} \\
    \midrule
    \texttt{CMakeLists.txt} &
      Removed 20 duplicate seismic targets ($115 \to 88$). \\[3pt]
    \texttt{NonlinearSubModelEvolver.hh} &
      Public \texttt{section\_tangent()}/\texttt{section\_forces()} aliases;
      public \texttt{commit\_state()}/\texttt{revert\_state()};
      factory injection; \texttt{static\_assert} for concept. \\[3pt]
    \texttt{header\_files.hh} &
      Added includes for all 6 new headers. \\
    \bottomrule
  \end{tabular}
\end{table}


% -------------------------------------------------------------------------
\subsection{Remaining Work}
\label{sec:refactor-remaining}

The following phases from the original roadmap
(Chapter~\ref{ch:multiscale-architecture}) are deferred to future
chapters:

\begin{description}
  \item[Phase 0.2 -- PCH:] Precompiled headers for third-party
    includes (Eigen, PETSc, VTK).
  \item[Phase 0.3 -- Explicit instantiation:] \texttt{extern template}
    for the most-used template combinations.
  \item[Phase 4.4 -- Main driver refactoring:]
    Replace the staggered loop in
    \texttt{main\_lshaped\_multiscale.cpp} with
    \texttt{MultiscaleAnalysis::step()}.
  \item[Phase 6 -- Parallelism:] OpenMP parallel sub-model solves
    and perturbation tangent.
  \item[Phase 7 -- SRP cleanup:] Extract VTK export and crack
    tracking from the evolver into separate classes.
  \item[Phase 8 -- Header modularisation:] Split
    \texttt{header\_files.hh} into tiered includes.
\end{description}
